\subsection{Network Architecture: PMDF-Net}
\label{sec:architecture}

PMDF-Net is a fully convolutional encoder-decoder network that jointly processes laser ultrasonic signals in the time domain and time-frequency domain. The architecture follows a dual-branch encoder design with a shared decoder, enabling complementary feature extraction across domains before reconstruction of the denoised time-domain waveform.

\textbf{Overall structure.} PMDF-Net consists of three stages: a time-domain encoder branch, a time-frequency encoder branch, a fusion module at the bottleneck, and a shared decoder. Both encoder branches progressively downsample their input representations to capture multi-scale features, then the decoder upsamples and fuses information via skip connections to reconstruct the denoised output. The network takes a preprocessed input signal $x$ and outputs a denoised waveform estimate $\hat{y}$.

\textbf{Time-domain branch.} The time-domain branch operates directly on the 1D input waveform. It comprises approximately four to five encoder blocks, each consisting of a 1D convolutional layer with a kernel size of 3 or 5 followed by batch normalization and a ReLU activation. Downsampling is performed via convolutional layers with stride 2, progressively halving the temporal resolution while increasing the channel dimension to capture features at increasingly large spatial scales. The decoder mirrors this structure with a corresponding number of upsampling blocks (transposed convolutions or strided upsampling) and skip connections to the encoder feature maps in the U-Net style \cite{ronneberger2015unet}. These skip connections preserve high-frequency waveform details that would otherwise be lost through successive downsampling. Rather than predicting the clean signal directly, the time-domain branch employs residual learning \cite{zhang2017denoising}: it predicts the noise residual $r$ such that $\hat{y} = x - r$. This formulation simplifies the learning objective by letting the network focus on the difference between noisy and clean signals, which has been shown to improve convergence and denoising performance in convolutional architectures.

\textbf{Time-frequency branch.} The time-frequency branch first transforms the input waveform into a time-frequency representation via STFT. Using the preprocessing parameters from Section~\ref{sec:preprocessing}, the STFT uses a Hanning window of length $W = 256$ with hop size $H = 64$, yielding a spectrogram with frequency resolution $\Delta f \approx 60$\,kHz and time resolution $\approx 0.8$\,$\mu$s. The branch then processes the STFT magnitude spectrogram---treated as a single-channel 2D image---with a 2D convolutional encoder-decoder. The encoder applies 2D convolutions with kernel sizes of $3 \times 3$ or $5 \times 5$, batch normalization, and ReLU, progressively downsampling the spectrogram to capture both low-frequency dispersive components and higher-frequency arrivals. The decoder upsamples via transposed convolutions and concatenates skip connections with the encoder, following the same U-Net topology used by the time-domain branch. After the encoder-decoder processing, an inverse STFT reconstructs the time-domain component, which is concatenated with features from the time-domain branch at the fusion module.

\textbf{Fusion module.} At the bottleneck, the feature representations from both branches are fused via concatenation along the channel dimension, yielding a combined feature tensor that encodes both time-domain waveform structure and time-frequency spectral characteristics. A $1 \times 1$ (pointwise) convolution reduces the concatenated channels to a compact representation, followed by a ReLU non-linearity to enable the fusion to learn non-trivial combinations of the two domains. This fused bottleneck representation is passed to the shared decoder.

\textbf{Shared decoder.} The decoder takes the fused bottleneck representation and progressively upsamples it to the original signal length via transposed convolutions. Skip connections from the corresponding encoder blocks inject high-resolution features at each scale, preserving fine temporal details in the output waveform. Because the network is fully convolutional, PMDF-Net accepts input signals of arbitrary length, provided the length is a multiple of the total downsampling factor (which is $2^4 = 16$ for a four-block architecture).

\textbf{Key design decisions.} The fully convolutional design ensures that the network can process variable-length signals without modification, making it suitable for real-time or batch processing of laser ultrasonic waveforms. Skip connections in both branches follow the U-Net principle of preserving high-frequency information through the encoding process. The absence of recurrent layers keeps inference fast and avoids gradient propagation challenges common in RNN-based denoisers. The total parameter count is kept in the range of approximately 1--5 million parameters, sufficiently large to learn complex noise patterns while remaining within practical memory and compute constraints for desktop GPU inference.

\textbf{Input and output.} The network accepts a preprocessed 1D signal $x$ of fixed length $L = 1024$ samples as the primary input. The time-frequency branch additionally requires the STFT of $x$ computed during preprocessing. The output is a denoised waveform estimate $\hat{y}$ of the same length as the input, from which the noise residual can be directly computed as $\hat{r} = x - \hat{y}$.